\expandafter\def\csname HistoricalNoteText26\endcsname{Genau genommen folgt aus der vorangehenden Analysis nur, dass $-G'$ und $-G''$ die Gleichung 13 erfüllen müssen; daraus scheint ein Zweifel zu entstehen, ob es nicht zulässig sei, beide, $-G'$ und $-G''$, \textit{derselben} negativen Wurzel gleichzusetzen und die dritte Wurzel völlig zu vernachlässigen. Man wird aber leicht erkennen, sofern die zweite und dritte Wurzel der Gleichung ungleich sind, dass aus $-G'=\allowbreak -G''$ die Beziehungen $\gamma'=\gamma''$, $\alpha'=\alpha''$, $\beta'=\beta''$ und folglich
$-\alpha'\alpha''-\beta'\beta''+\allowbreak \gamma'\gamma''=\allowbreak -\alpha'\alpha'-\beta'\beta'+\allowbreak \gamma'\gamma'=1$ folgen, was der vierten Gleichung in [1] widerspricht. Vergleiche, was unten über den Fall zweier gleicher Wurzeln der Gleichung 13 gesagt werden wird.}
